% !PRINTED 20
\weberpage{20}
%<<CONT
Die Gleichheit findet hier nur dann statt, wenn der Quotient $z:c$
reell und negativ ist.

\begin{center}
\rule{0.18\linewidth}{0.4pt}
\end{center}

Es ist noch ein Wort über das wichtigste Hülfsmittel der Algebra, die
Buchstabenrechnung, zu sagen. Die Anwendung dieses Hülfsmittels ist so
allgemein, dass man bisweilen das Wort Buchstabenrechnung geradezu
synonym mit Algebra gebraucht. Die Regeln, wie mit solchen
Buchstabenausdrücken gerechnet wird, setzen wir als bekannt voraus.
Gleichungen zwischen Buchstabenausdrücken können von zweierlei Art
sein; entweder es sind sogenannte Identitäten, d.~h. die zwei einander
gleich gesetzten Ausdrücke können durch Anwendung der Rechenregeln so
umgeformt werden, dass beide Ausdrücke genau übereinstimmen. Man erhält
dann aus solchen Buchstabengleichungen richtige Zahlengleichungen, wenn
die Buchstaben durch irgend welche, sei es reelle, sei es complexe
Zahlen ersetzt werden, vorausgesetzt, dass dabei nicht die Forderung
der Division durch Null auftritt. Die Buchstaben in solchen Gleichungen
werden oft auch als Variable bezeichnet, weil man sich vorstellen
kann, ohne je zu einem Widerspruch zu gelangen, dass für die Buchstaben
nach und nach andere und andere Zahlwerthe gesetzt werden.

Eine andere Art von Gleichungen zwischen Buchstabenausdrücken haben
nicht diesen Charakter der Identität. Sie enthalten vielmehr eine
Forderung, die an die Zahlen gestellt werden, die man ohne einen Fehler
zu begehen für die Buchstaben einsetzen darf. Die Algebra hat die
Aufgabe, Zahlwerthe zu ermitteln, die einer solchen Forderung genügen,
die Gleichung zu lösen. In diesen Gleichungen werden die Buchstaben
auch als \gesperrt{„Unbekannte“} bezeichnet. Es kommen sehr häufig in
ein und derselben Gleichung Buchstaben von zwei Arten vor, solche, für
die beliebige Zahlwerthe gesetzt werden sollen, und andere, deren
Zahlwerth erst ermittelt werden soll.

\begin{center}
\rule{0.24\linewidth}{0.4pt}
\end{center}
